\chapter{Planificación}

\section{Introducción}

En este capítulo se explica la metodología y ritmo de trabajo que se seguirá durante el desarrollo del proyecto. Con respecto a la metodología, se empleará un modelo personalizado basado en distintos aspectos de la metodología ágil Scrum. Cabe destacar que, pese a que se usarán elementos de Scrum, no se seguirá esta metodología, debído a que está orientada al trabajo en equipo y, por tanto, no se adapta completamente a un proyecto individual como este. \\

Las tareas se organizarán en historias de usuario que se desarrollarán en sprints de dos semanas de duración. Al final de cada sprint, se tendrá una reunión con el tutor del TFG en la que se revisará el trabajo realizado, se identificarán posibles mejoras y se marcarán los siguientes objetivos a abordar.

\section{Cálculo de la velocidad}

\section{Historias de usuario}
A continuación se expone un listado con todas las historias de usuario que contiene la aplicación:

\begin{itemize}
    \item Código HU1: Como cliente quiero poder crearme un perfil para poder usar la aplicación.
    \item Código HU2: Como cliente quiero poder editar mi perfil para modificar datos erroneos.
    \item Código HU3: Como cliente quiero poder eliminar mi perfil para eliminar mis datos de la aplicación.
    \item Código HU4: Como cliente quiero poder iniciar sesión en la aplicación para poder tomar decisiones.
    \item Código HU5: Como cliente quiero poder cerrar sesión en la aplicación para que nadie más pueda acceder a mi cuenta.
    \item Código HU6: Como cliente quiero poder crear un grupo para tomar decisiones con mis amigos.
    \item Código HU7: Como cliente quiero poder salirme de un grupo para dejar de formar parte de él.
    \item Código HU8: Como cliente quiero poder añadir miembros a un grupo para añadirlos en la toma de decisiones.
    \item Código HU9: Como cliente quiero poder eliminar miembros de un grupo para que dejen de formar parte del grupo
    \item Código HU10: Como cliente quiero poder iniciar una votación para tomar una decisión.
    \item Código HU11: Como cliente quiero asignarle un tiempo máximo a una votación para recibir los resultados en tiempo limitado.
    \item Código HU12: Como cliente quiero poder votar en las votaciones para tener derecho de elección.
    \item Código HU13: Como cliente quiero poder añadir opciones a una votación para poder votarlas posteriormente.
    \item Código HU14: Como cliente quiero poder elegir el tipo de votación para adaptarlo a mis necesidades.
    \item Código HU15: Como cliente quiero ver las decisiones pasadas para recordar las elecciones tomadas.
    \item Código HU16: Como cliente quiero ver la opción ganadora en una votación para saber qué decisión tomar. 
    \item Código HU17: Como cliente quiero poder seleccionar decisiones predeterminadas para facilitar mi toma de decisiones.
    \item Código HU18: Como cliente quiero poder subir imágenes a la aplicación para usarlas en las opciones o como foto de perfil o grupo.
    \item Código HU19: Como administrador quiero poder crear otro administrador para que me ayude en las tareas de gestión.
    \item Código HU20: Como administrador quiero poder eliminar a otro administrador para que no pueda hacer cambios en la aplicación.
    \item Código HU21: Como administrador quiero poder iniciar sesión en la aplicación para hacer tareas de gestión.
    \item Código HU22: Como administrador quiero poder cerrar sesión en la aplicación para que nadie más pueda acceder a mi cuenta.
    \item Código HU23: Como administrador quiero poder modificar los datos de un usuario para solucionar problemas de gestión.
    \item Código HU24: Como administrador quiero poder modificar datos de un grupo para solucionar problemas de gestión.
    \item Código HU25: Como administrador quiero poder modificar datos de una votación para solucionar problemas de gestión.
    \item Código HU26: Como administrador quiero poder eliminar un usuario para solucionar problemas de gestión.
    \item Código HU27: Como administrador quiero poder eliminar un grupo para solucionar problemas de gestión.
    \item Código HU28: Como administrador quiero poder eliminar una votación para solucionar problemas de gestión.
    \item Código HU29: Como administrador quiero poder realizar las tareas de un cliente para cercionarme de errores.
\end{itemize}    


\section{Planificación temporal}

\section{Diagrama de Gantt}


\section{Estimación de costes}